\section{Solution Architecture and Model Selection}
\label{sec:model_selection}

Two architectural paradigms were developed and empirically evaluated under identical experimental conditions across all 15 MVTec AD categories: the non-parametric PatchCore feature-matching framework and a deep generative convolutional autoencoder (the Keras CAE). Both pipelines were trained and calibrated under the deterministic 85\% fitting / 15\% normal validation split with seed 42, evaluated at the canonical $256 \times 256$ spatial resolution, and scored against unseen evaluation splits under our strict zero-test-leakage protocol described in Section~\ref{sec:threshold_calibration}.

\subsection{PatchCore: Deep Feature Transfer and Coreset Memory Banks}

The production-grade engine recommended for plant-wide deployment is PatchCore utilising a frozen ResNet-18 backbone pre-trained on ImageNet. Beyond superior empirical detection metrics, this architecture exhibits three decisive practical advantages that resolve longstanding operational pain points in manufacturing IT:

First, the architecture eliminates complex model training on the factory floor. The core neural network (ResNet-18) acts as a fixed feature extractor. Once deployed, the factory computer never needs to train models or tune parameters. Adding a new product type takes less than fifteen seconds, requiring only a single forward pass over 100--200 standard reference images to build the baseline.

Second, the memory bank guarantees complete client data sovereignty and protection of intellectual property. Conforming patterns are stored locally in a reference memory bank rather than in network weights. Because this representation does not alter the underlying model, proprietary part geometries, CAD toolpaths, and confidential material formulations remain strictly confined to the local industrial personal computer (IPC) on the factory floor; raw production images are never transmitted to external cloud servers. To guarantee fast inspection times that easily match conveyor speeds, a coreset selection algorithm compresses the stored reference library by an order of magnitude whilst retaining the essential patterns of conforming components.

Third, the architecture avoids expensive specialised hardware through deliberate model selection. Whilst heavier neural networks (such as Wide-ResNet-50-2) were evaluated during exploratory research, their large memory footprint requires roughly 11.8\,GB of RAM, overloading standard industrial PCs. Standardising on ResNet-18 reduces memory usage to under 3.8\,GB. This represents a deliberate systems engineering choice: sacrificing less than one percentage point of academic $F_1$ score to eliminate the need for costly industrial graphics cards (+€5,000 to €8,000 per inspection station) and enable reliable 65--95\,ms inspection on standard, fanless industrial PCs using standard processor instructions.

\subsection{Comparative Analysis: The Failure Modes of Generative Autoencoders}

The alternative paradigm evaluated is a fully convolutional autoencoder implemented in Keras, trained to reconstruct nominal images via a composite loss function combining the Structural Similarity Index Measure (SSIM) and Mean Squared Error (MSE), augmented by masked image modelling (MIM). During inference, the pixel-wise reconstruction residual serves as the continuous anomaly map, from which image-level classification scores are extracted via spatial aggregation.

Whilst generative reconstruction models offer conceptual appeal -- learning a compact latent distribution of nominal appearance and quantifying defects via reconstruction failure -- they suffer from a fatal structural flaw in industrial practice. Pixel-level reconstruction objectives inherently minimise global image variance rather than local defect contrast. On high-frequency stochastic textures -- such as \emph{Wood} grain, \emph{Carpet} pile, or \emph{Leather} surfaces -- the generative decoder yields slightly smoothed reconstructions of high-frequency spatial patterns. This slight reconstruction blur generates elevated residual errors across conforming surfaces, inducing severe pseudo-scrap and artificially deflating the operational $F_1$-score. Moreover, the resulting spatial heatmaps diffuse across the entire component surface rather than isolating localised defect boundaries. Table~\ref{tab:paradigm_comparison} quantifies this performance disparity across all 15 benchmark categories.

\begin{table}[bp]
\centering
\caption{Macro-averaged performance across all 15 MVTec AD categories under our strict zero-test-leakage protocol. PatchCore ResNet-18 constitutes the production-grade reference configuration.}
\label{tab:paradigm_comparison}
\footnotesize
\resizebox{\linewidth}{!}{%
\begin{tabular}{@{}lccccc@{}}
\toprule
\textbf{Architecture} & \textbf{Image $F_1$ $\uparrow$} & \textbf{Image PR-AUC $\uparrow$} & \textbf{AUPIMO $\uparrow$} & \textbf{Pixel AUROC $\uparrow$} & \textbf{Client Training} \\
\midrule
PatchCore (ResNet-18)       & \textbf{0.929} & \textbf{0.988} & \textbf{0.603} & 0.968 & None ($<$15\,s forward pass) \\
Keras CAE (MIM + SSIM+MSE)  & 0.495          & 0.867          & 0.058          & 0.875 & Iterative GPU training \\
\bottomrule
\end{tabular}%
}
\end{table}

Empirically, the Keras CAE trails the PatchCore architecture by an order of magnitude on pixel-level AUPIMO ($0.058$ versus $0.603$) and by $0.434$ points in image $F_1$-score ($0.495$ versus $0.929$). This deficit cannot be mitigated through hyperparameter optimisation or alternative loss weighting; it reflects an intrinsic limitation of pixel-space generative modelling for industrial surface inspection.

Furthermore, training of the Keras CAE was executed using the AdamW optimiser with an initial learning rate $\eta_0 = 1 \times 10^{-4}$ and weight decay of $1 \times 10^{-4}$, capped at a maximum budget of 100 epochs with dynamic early stopping (patience 20 epochs) and learning-rate halving on validation plateaus (patience 5 epochs). It necessitated a restricted mini-batch size of 16 to comply with the physical memory bounds of edge hardware. In deep generative networks, small mini-batch regimes introduce high gradient variance and destabilise batch normalisation statistics, impeding smooth convergence onto the nominal manifold. Whilst scaling batch sizes to 32 or 64 could marginally stabilise training dynamics, doing so would require server-grade GPU infrastructure, directly violating the core requirement of localised, low-overhead factory deployment.

\subsection{Continuous Calibration: Mitigating Operational Drift Without Retraining}

A decisive operational advantage of the decoupled coreset architecture is its capability to absorb manufacturing distribution drift without requiring model retraining. When minor physical alterations occur on the line -- such as a new steel batch exhibiting marginal surface roughness variations, a slight dye shade shift from an upstream textile supplier, or illumination drift caused by lamp aging -- conventional deep learning pipelines begin generating elevated anomaly scores for nominal parts. In a standard convolutional network, resolving this issue requires collecting fresh image batches, executing full backpropagation, re-validating convergence, and redeploying weights -- a procedure demanding hours or days of specialised engineering intervention.

Within the PatchCore framework, this drift is resolved deterministically at line side. Quality assurance operators examine the flagged units via the local inspection terminal and confirm them as false positives. The system extracts feature patch embeddings from the confirmed nominal components and appends them directly to the active coreset memory bank. An incremental coreset subsampling step restores optimal memory density within seconds, and the decision threshold is updated against the expanded normal validation pool. The network parameters remain completely untouched. This enables line-side personnel to adapt the inspection system to legitimate process drift within thirty seconds, requiring zero data science expertise.
